%!TEX root = ../main.tex
% =============================================================================
%  参考文献（格式规范第七条：引用他人或公开资料成果必须列出参考文献，
%            并在正文引用处予以标注；格式按 GB/T 7714）
%
%  ⚠️ 三条引用纪律
%     1 宁可少引三条，不可引错一条——引用前必须核对官方原文
%       （国家发展改革委 / 国家能源局官网、\url{https://openstd.samr.gov.cn/）}
%     2 不得编造政策文号、标准号、企业数据、参考文献（属学术不端）
%     3 行业数字只能用于背景与合理性检验，不得当作本题计算结果，
%       不得出现在结果表或「计算结果」章节
%
%  本节只保留三条经官网原文核对通过的政策文件；文号、发文机关与成文日期逐条核对，
%  无法核实的条目一律删除，不保留占位。
% =============================================================================

% thebibliography 环境自带标题（ctexart 的 \refname 为「参考文献」），
% 因此不再手写 \section*{参考文献}，避免出现两个标题。
\renewcommand{\refname}{参考文献}

\begin{thebibliography}{99}
  \setlength{\itemsep}{2pt}

  \bibitem{ref136}
  国家发展改革委, 国家能源局.
  关于深化新能源上网电价市场化改革~促进新能源高质量发展的通知:
  发改价格〔2025〕136号[Z]. 2025-01-27.
  \url{https://www.gov.cn/zhengce/zhengceku/202502/content_7002959.htm}
  \note{正文第一章引用；发文机关、文号与成文日期按中国政府网「国务院部门文件」栏目核对（2026-09-10）}

  \bibitem{ref1093}
  国家发展改革委.
  关于进一步完善分时电价机制的通知:
  发改价格〔2021〕1093号[Z]. 2021-07-26.
  \url{https://www.gov.cn/zhengce/zhengceku/2021-07/29/content_5628297.htm}
  \note{分时电价的时段划分与峰谷价差机制出处；本文用于分时电价骨架与峰谷价比的背景说明}

  \bibitem{ref1217}
  国家发展改革委, 国家能源局.
  电力现货市场基本规则（试行）:
  发改能源规〔2023〕1217号[Z]. 2023-09-07.
  \url{https://www.gov.cn/zhengce/zhengceku/202309/content_6904881.htm}
  \note{「日前计划与实时偏差」结算机制的政策原型；本文问题 2、3 的信息结构与偏差定价口径参考此机制}

\end{thebibliography}
